\documentclass[12pt,spanish]{article}
\usepackage[utf8]{inputenc}
\usepackage[spanish]{babel}
\usepackage{amsmath, amssymb, amsthm}
\usepackage{graphicx}
\usepackage{hyperref}
\usepackage[top=2.5cm, bottom=2.5cm, left=3cm, right=2.5cm]{geometry}
\usepackage{float}
\usepackage{longtable}
\usepackage{array}
\usepackage{booktabs}
\usepackage{xcolor}
\usepackage{enumitem}
\usepackage{titlesec}
\usepackage{fancyhdr}
\usepackage{multicol}
\usepackage{caption}
\usepackage[most]{tcolorbox}

% Configuración de página
\pagestyle{fancy}
\fancyhf{}
\rhead{\thepage}
\lhead{\textit{Proyecto de Investigación Pedagógica - Lengua Castellana}}
\renewcommand{\headrulewidth}{0.4pt}

% Configuración de títulos
\titleformat{\section}{\Large\bfseries}{\thesection}{1em}{}
\titleformat{\subsection}{\large\bfseries}{\thesubsection}{1em}{}
\titleformat{\subsubsection}{\normalsize\bfseries}{\thesubsubsection}{1em}{}

% Colores
\definecolor{color1}{RGB}{0,102,204}
\definecolor{color2}{RGB}{204,102,0}
\definecolor{color3}{RGB}{0,153,76}

% Comandos personalizados
\newcommand{\programa}[1]{\texttt{\color{color2}#1}}
\newcommand{\autor}{Hernank10}

% Cajas para notas
\newtcolorbox{nota}{
    colback=blue!5,
    colframe=blue!50!black,
    arc=4pt,
    boxrule=1pt,
    title=Nota pedagógica
}

\newtcolorbox{ejemplo}{
    colback=green!5,
    colframe=green!50!black,
    arc=4pt,
    boxrule=1pt,
    title=Ejemplo en el aula
}

\newtcolorbox{codigo}{
    colback=gray!10,
    colframe=gray!50!black,
    arc=4pt,
    boxrule=1pt,
    title=Código / Comando
}

% Documento
\begin{document}

% Portada
\begin{titlepage}
    \centering
    \vspace*{2cm}
    {\Huge \textbf{Impacto de un Ecosistema Digital Interactivo en el Desarrollo de Competencias de Lectura, Escritura y Sintaxis en la Enseñanza de la Lengua Castellana} \par}
    \vspace{1cm}
    {\Large Proyecto de Investigación Pedagógica \par}
    \vspace{1.5cm}
    {\large \textbf{Autor:} \autor \par}
    \vspace{0.5cm}
    {\large \textbf{Área:} Didáctica de la Lengua Castellana \par}
    \vspace{0.5cm}
    {\large \textbf{Duración:} 8 meses \par}
    \vspace{0.5cm}
    {\large \textbf{Nivel educativo:} Educación Secundaria (14-18 años) \par}
    \vspace{2cm}
    {\large \today \par}
    \vfill
    {\large \textbf{Repositorio:} \url{https://github.com/Hernank10/mis-programas-python} \par}
    \vspace{0.3cm}
    {\large \textbf{Sitio web:} \url{https://hernank10.github.io/mini-htmls-educativos/} \par}
\end{titlepage}

% Resumen
\section*{Resumen}
\addcontentsline{toc}{section}{Resumen}

Este proyecto de investigación pedagógica evalúa el impacto de un ecosistema digital interactivo basado en Python y HTML educativo en el desarrollo de competencias lingüísticas en lengua castellana. Mediante un diseño cuasi-experimental con grupo control y grupo experimental (N=120 estudiantes), se medirán las competencias de lectura crítica, producción textual, dominio sintáctico y ortográfico durante 12 semanas de intervención. Los recursos incluyen más de 80 mini lecciones HTML interactivas, aplicaciones Python, juegos gramaticales y talleres prácticos. Se hipotetiza que el grupo experimental mostrará una mejora significativa (≥25\%) en las pruebas estandarizadas en comparación con la enseñanza tradicional.

\textbf{Palabras clave:} Lengua castellana, innovación pedagógica, tecnología educativa, lectura crítica, producción textual, sintaxis, HTML educativo, Python.

\newpage

% Tabla de contenido
\tableofcontents
\newpage

\section{Planteamiento del Problema}

\subsection{Diagnóstico actual de la enseñanza de la lengua castellana}

Las investigaciones en educación lingüística han identificado varias problemáticas persistentes, resumidas en la Tabla \ref{tab:problemas}.

\begin{table}[H]
\centering
\caption{Problemáticas en la enseñanza de la lengua castellana}
\label{tab:problemas}
\begin{tabular}{|p{4cm}|p{6cm}|p{4cm}|}
\hline
\textbf{Problema} & \textbf{Evidencia empírica} & \textbf{Consecuencia pedagógica} \\
\hline
Baja motivación lectora & Solo el 45\% de estudiantes lee por placer (PISA 2022) & Dificultad para comprender textos complejos \\
\hline
Déficit en producción textual & 6 de cada 10 estudiantes no alcanza nivel mínimo en escritura & Limitación para expresar ideas coherentemente \\
\hline
Errores sintácticos recurrentes & Confusión entre sujeto/predicado, dequeísmo, queísmo & Pobreza estructural en oraciones \\
\hline
Falta de retroalimentación inmediata & Correcciones con 5-7 días de retraso & Errores fosilizados \\
\hline
Unidimensionalidad evaluativa & Solo exámenes escritos & No se evalúa proceso ni diversidad de estilos \\
\hline
\end{tabular}
\end{table}

\subsection{Pregunta de investigación}

\begin{quote}
\textbf{¿Cómo influye el uso del ecosistema digital interactivo (mini-htmls-educativos + programas Python) en el desarrollo de competencias de lectura crítica, producción textual y dominio sintáctico en estudiantes de secundaria?}
\end{quote}

\subsection{Preguntas específicas}

\begin{enumerate}
    \item ¿Existe una mejora significativa en la comprensión lectora después de la intervención con mini lecciones interactivas?
    \item ¿Qué efecto tiene la retroalimentación inmediata en la reducción de errores sintácticos (sujeto-predicado, dequeísmo)?
    \item ¿Cómo influye la gamificación (juegos gramaticales) en la frecuencia y calidad de la producción escrita voluntaria?
    \item ¿Qué relación existe entre el número de lecciones completadas y la mejora en pruebas estandarizadas de lengua?
\end{enumerate}

\subsection{Hipótesis}

\begin{table}[H]
\centering
\caption{Hipótesis de investigación}
\label{tab:hipotesis}
\begin{tabular}{|p{3cm}|p{7cm}|p{4cm}|}
\hline
\textbf{Hipótesis} & \textbf{Enunciado} & \textbf{Indicador} \\
\hline
H1 & El grupo experimental mejorará significativamente (p < 0.05, d > 0.5) en comprensión lectora & Diferencia de medias en prueba estandarizada \\
\hline
H2 & La tasa de errores sintácticos se reducirá en al menos un 35\% en el grupo experimental & Análisis de textos producidos \\
\hline
H3 & Los estudiantes dedicarán ≥ 1.5 horas semanales de práctica voluntaria & Registros automáticos del LMS \\
\hline
H4 & El 80\% de los estudiantes calificará la experiencia como útil o muy útil (≥ 4/5) & Encuesta de percepción \\
\hline
\end{tabular}
\end{table}

\newpage
\section{Marco Teórico}

\subsection{Fundamentos pedagógicos y lingüísticos}

La Tabla \ref{tab:teorias} presenta las bases teóricas que sustentan este proyecto.

\begin{table}[H]
\centering
\caption{Fundamentos pedagógicos del proyecto}
\label{tab:teorias}
\begin{tabular}{|p{3.5cm}|p{3.5cm}|p{6cm}|}
\hline
\textbf{Teoría/Enfoque} & \textbf{Autor(es)} & \textbf{Aplicación en el proyecto} \\
\hline
Constructivismo social & Vygotsky (1978) & Zona de Desarrollo Próximo: las lecciones se adaptan al nivel del estudiante \\
\hline
Aprendizaje activo & Bonwell \& Eison (1991) & Ejercicios interactivos donde el estudiante escribe, clasifica y transforma \\
\hline
Enfoque comunicativo & Hymes (1972), Canale (1983) & Lecciones en contextos reales (redacción médica, científica, web) \\
\hline
Gamificación & Kapp (2012) & Juegos de sintaxis con puntuaciones y niveles \\
\hline
Retroalimentación correctiva & Hattie \& Timperley (2007) & Efecto tamaño > 0.7: la inmediatez potencia el aprendizaje \\
\hline
Lingüística textual & Van Dijk (1980), Cassany (2006) & Énfasis en producción de textos completos, no solo oraciones \\
\hline
\end{tabular}
\end{table}

\subsection{Competencias lingüísticas evaluadas}

Basado en el \textbf{Marco Común Europeo de Referencia (MCER)} y los estándares de la \textbf{RAE}:

\begin{itemize}
    \item \textbf{Comprensión lectora:} Extraer información, inferir y reflexionar
    \item \textbf{Producción textual:} Escribir textos coherentes, cohesionados y adecuados
    \item \textbf{Dominio sintáctico:} Construir oraciones correctas y variadas
    \item \textbf{Competencia ortográfica:} Usar correctamente tildes, diptongos, hiatos
    \item \textbf{Competencia discursiva:} Organizar textos según propósito y audiencia
\end{itemize}

\subsection{Estado del arte}

\begin{table}[H]
\centering
\caption{Investigaciones previas relevantes}
\label{tab:estado_arte}
\begin{tabular}{|p{3.5cm}|p{5cm}|p{5cm}|}
\hline
\textbf{Estudio} & \textbf{Aporte} & \textbf{Limitación} \\
\hline
Cassany (2006) & Importancia de la escritura digital & Poca interactividad \\
\hline
Marquès (2010) & Software educativo mejora motivación & Contenido cerrado \\
\hline
Area \& Pessoa (2012) & Alfabetización digital necesaria en aula & Falta de recursos específicos \\
\hline
\textbf{Este proyecto} & Ecosistema unificado y gratuito & - \\
\hline
\end{tabular}
\end{table}

\newpage
\section{Objetivos}

\subsection{Objetivo general}

Evaluar el impacto de un ecosistema digital interactivo basado en HTML y Python en el desarrollo de competencias de lectura crítica, producción textual y dominio sintáctico en estudiantes de secundaria.

\subsection{Objetivos específicos}

\begin{table}[H]
\centering
\caption{Objetivos específicos e instrumentos de medición}
\label{tab:objetivos}
\begin{tabular}{|p{6cm}|p{7cm}|}
\hline
\textbf{Objetivo} & \textbf{Instrumento de medición} \\
\hline
Medir mejora en comprensión lectora & Pre-test / post-test (20 ítems) \\
\hline
Evaluar reducción de errores sintácticos & Análisis de textos producidos \\
\hline
Cuantificar relación entre práctica voluntaria y mejora & Análisis de correlación (Pearson) \\
\hline
Documentar percepción de usabilidad y motivación & Encuesta de escala Likert + grupo focal \\
\hline
\end{tabular}
\end{table}

\newpage
\section{Metodología}

\subsection{Diseño de investigación}

\begin{itemize}
    \item \textbf{Tipo:} Cuasi-experimental con grupo control no equivalente, con mediciones pre-test y post-test, y seguimiento a 30 días.
    \item \textbf{Duración:} 12 semanas (un trimestre académico)
\end{itemize}

\textbf{Esquema:}

\begin{center}
\fbox{\parbox{0.9\textwidth}{
G1 (Control): O1 → Xtradicional → O2 → O3\\
G2 (Experimental): O1 → XDigital (12 semanas) → O2 → O3
}}
\end{center}

\subsection{Participantes}

\begin{table}[H]
\centering
\caption{Características de los participantes}
\label{tab:participantes}
\begin{tabular}{|p{4cm}|p{8cm}|}
\hline
\textbf{Característica} & \textbf{Detalle} \\
\hline
N & 120 estudiantes (60 control, 60 experimental) \\
\hline
Edad & 14-16 años (3º y 4º de ESO) \\
\hline
Centros & 2 institutos públicos (60 estudiantes cada uno) \\
\hline
Nivel socioeconómico & Medio-bajo \\
\hline
Horas de lengua/semana & 4 horas \\
\hline
Metodología control & Libro de texto + ejercicios en papel \\
\hline
Metodología experimental & Mini lecciones HTML + programas Python \\
\hline
\end{tabular}
\end{table}

\subsection{Instrumentos de recolección}

\begin{table}[H]
\centering
\caption{Instrumentos de recolección de datos}
\label{tab:instrumentos}
\begin{tabular}{|p{4cm}|p{4cm}|p{5cm}|}
\hline
\textbf{Instrumento} & \textbf{Aplicación} & \textbf{Qué mide} \\
\hline
Prueba de comprensión lectora (pre/post) & Semana 1 y 12 & Niveles literal, inferencial, crítico \\
\hline
Análisis de errores sintácticos & Semanas 4, 8, 12 & Concordancia, dequeísmo, puntuación \\
\hline
Registro automático & Semanas 1-12 & Práctica voluntaria \\
\hline
Encuesta de percepción & Semana 12 & Usabilidad, motivación, utilidad \\
\hline
Grupo focal & Semana 13 & Experiencia cualitativa \\
\hline
\end{tabular}
\end{table}

\subsection{Programa de intervención (12 semanas)}

\begin{longtable}{|p{2cm}|p{4cm}|p{6cm}|}
\hline
\textbf{Semana} & \textbf{Tema} & \textbf{Recurso digital} \\
\hline
1 & Sujeto y predicado & \programa{sujeto\_predicado.html} \\
\hline
2 & Tipos de párrafos & \programa{tipos-parrafo.html} \\
\hline
3 & Conectores gramaticales & \programa{conectores.html} \\
\hline
4 & Oraciones impersonales & \programa{oraciones-impersonales.py} \\
\hline
5 & Diptongos e hiatos & \programa{diptongos\_hiatos.html} \\
\hline
6 & Signos de puntuación & \programa{signos-puntuacion.html} \\
\hline
7 & Técnicas de redacción & \programa{tecnicas-redaccion.html} \\
\hline
8 & Redacción académica & \programa{redaccion-licenciatura.html} \\
\hline
9 & Juego de sintaxis & \programa{juego-sintaxis.py} \\
\hline
10 & Redacción especializada & \programa{redaccion-medicina-colombia.html} \\
\hline
11 & Repaso general & Menú principal + elección libre \\
\hline
12 & Evaluación final & Post-test + encuesta \\
\hline
\caption{Secuencia didáctica semanal}
\label{tab:programas}
\end{longtable}

\newpage
\section{Análisis de Datos}

\subsection{Variables}

\begin{table}[H]
\centering
\caption{Variables del estudio}
\label{tab:variables}
\begin{tabular}{|p{3cm}|p{4cm}|p{6cm}|}
\hline
\textbf{Tipo} & \textbf{Variable} & \textbf{Indicador} \\
\hline
Independiente & Metodología & Tradicional (0) / Digital (1) \\
\hline
Dependiente principal & Comprensión lectora & Puntuación en post-test (0-100) \\
\hline
Dependientes secundarias & Producción textual & Rúbrica (4 dimensiones) \\
\hline
 & Dominio sintáctico & Porcentaje de aciertos \\
\hline
 & Tiempo de práctica & Registro automático \\
\hline
Control & Edad, género, nota previa & Covariables en ANCOVA \\
\hline
\end{tabular}
\end{table}

\subsection{Técnicas estadísticas}

\begin{table}[H]
\centering
\caption{Técnicas estadísticas a emplear}
\label{tab:estadisticas}
\begin{tabular}{|p{4cm}|p{4cm}|p{5cm}|}
\hline
\textbf{Contraste} & \textbf{Prueba} & \textbf{Propósito} \\
\hline
Normalidad & Shapiro-Wilk & Verificar distribución de datos \\
\hline
Diferencia intra-grupo & t de Student (muestras relacionadas) & Mejora dentro de cada grupo \\
\hline
Diferencia entre grupos & t de Student (muestras independientes) & Experimental vs control \\
\hline
Tamaño del efecto & d de Cohen & Magnitud de la diferencia \\
\hline
Relación práctica-mejora & Correlación de Pearson & Asociación lineal \\
\hline
Predictores de éxito & Regresión lineal múltiple & Modelo explicativo \\
\hline
Fiabilidad de encuesta & Alpha de Cronbach & Consistencia interna \\
\hline
\end{tabular}
\end{table}

\newpage
\section{Cronograma de Actividades}

\begin{table}[H]
\centering
\caption{Cronograma de actividades por mes}
\label{tab:cronograma}
\begin{tabular}{|p{3.5cm}|p{4cm}|p{3.5cm}|p{3.5cm}|}
\hline
\textbf{Mes} & \textbf{Actividad} & \textbf{Responsable} & \textbf{Entregable} \\
\hline
Mes 1 & Revisión bibliográfica + diseño instrumentos & Investigador & Marco teórico \\
\hline
Mes 2 & Capacitación docente (3 sesiones) & Investigador + profesorado & Guía de uso \\
\hline
Mes 2 & Solicitud de permisos éticos & Investigador & Cartas firmadas \\
\hline
Mes 3 & Aplicación de pre-test & Profesores & Datos de línea base \\
\hline
Meses 3-6 & Intervención (12 semanas) & Profesores & Registros automáticos \\
\hline
Mes 6 & Post-test + encuesta & Investigador & Datos post-intervención \\
\hline
Mes 7 & Prueba de retención & Investigador & Datos de memoria \\
\hline
Meses 7-8 & Análisis de datos & Investigador & Tablas, gráficos \\
\hline
Mes 8 & Redacción de informe final & Investigador & Informe + artículo \\
\hline
\end{tabular}
\end{table}

\subsection{Diagrama de Gantt simplificado}

\begin{center}
\begin{tabular}{|l|c|c|c|c|c|c|c|c|}
\hline
\textbf{Actividad} & \textbf{M1} & \textbf{M2} & \textbf{M3} & \textbf{M4} & \textbf{M5} & \textbf{M6} & \textbf{M7} & \textbf{M8} \\
\hline
Revisión bibliográfica & X & X & & & & & & \\
\hline
Diseño de instrumentos & & X & & & & & & \\
\hline
Permisos éticos & & X & & & & & & \\
\hline
Pre-test & & & X & & & & & \\
\hline
Intervención & & & X & X & X & X & & \\
\hline
Post-test & & & & & & X & & \\
\hline
Prueba de retención & & & & & & & X & \\
\hline
Análisis de datos & & & & & & & X & X \\
\hline
Redacción final & & & & & & & & X \\
\hline
\end{tabular}
\end{center}

\newpage
\section{Resultados Esperados}

\begin{table}[H]
\centering
\caption{Hipótesis y criterios de éxito}
\label{tab:resultados}
\begin{tabular}{|p{3cm}|p{5cm}|p{5cm}|}
\hline
\textbf{Hipótesis} & \textbf{Criterio cuantitativo} & \textbf{Criterio cualitativo} \\
\hline
H1 (comprensión lectora) & Diferencia ≥ 20\%, p < 0.05, d > 0.5 & Mejora en identificación de ideas principales \\
\hline
H2 (errores sintácticos) & Reducción ≥ 35\% de errores & Menos de 2 errores por texto \\
\hline
H3 (práctica voluntaria) & Media ≥ 1.5 horas/semana & Testimonios positivos \\
\hline
H4 (percepción) & ≥ 80\% puntúa utilidad ≥ 4/5 & Comentarios favorables \\
\hline
\end{tabular}
\end{table}

\subsection{Productos esperados}

\begin{enumerate}
    \item Informe final de investigación (80-100 páginas)
    \item Artículo científico para revista indexada
    \item Guía didáctica para docentes (40 páginas)
    \item Conjunto de datos anonimizados en repositorio abierto
    \item Video-resumen de 5 minutos para difusión
\end{enumerate}

\newpage
\section{Limitaciones y Estrategias de Mitigación}

\begin{table}[H]
\centering
\caption{Limitaciones y estrategias de mitigación}
\label{tab:limitaciones}
\begin{tabular}{|p{5cm}|p{8cm}|}
\hline
\textbf{Limitación} & \textbf{Estrategia de mitigación} \\
\hline
Efecto Hawthorne & Ambos grupos reciben atención similar; registros automáticos discretos \\
\hline
Diferencia basal entre centros & Centro como covariable en ANCOVA \\
\hline
Curva de aprendizaje tecnológico & Sesión de familiarización de 60 minutos \\
\hline
Fatiga por pantalla & Sesiones máximas de 50 minutos con pausas activas \\
\hline
Mortalidad muestral & Sobremuestreo inicial (N esperado 140) \\
\hline
Generalización externa & Descripción detallada del contexto; invitación a réplicas \\
\hline
\end{tabular}
\end{table}

\newpage
\section{Consideraciones Éticas}

\begin{table}[H]
\centering
\caption{Principios éticos y acciones correspondientes}
\label{tab:etica}
\begin{tabular}{|p{4cm}|p{9cm}|}
\hline
\textbf{Principio} & \textbf{Acción} \\
\hline
Consentimiento informado & Padres/tutores firman autorización \\
\hline
Confidencialidad & Datos anonimizados (códigos, no nombres) \\
\hline
Voluntariedad & Los estudiantes pueden retirarse sin penalización \\
\hline
Beneficencia & Recursos disponibles para todos los centros después del estudio \\
\hline
No maleficencia & Máximo 50 min/semana adicionales \\
\hline
\end{tabular}
\end{table}

\newpage
\section{Bibliografía}

\begin{enumerate}
    \item Area, M., \& Pessoa, T. (2012). De lo sólido a lo líquido: las nuevas alfabetizaciones ante los cambios culturales de la Web 2.0. \textit{Comunicar}, 38, 13-20.

    \item Bonwell, C. C., \& Eison, J. A. (1991). \textit{Active Learning: Creating Excitement in the Classroom}. ASHE-ERIC.

    \item Canale, M. (1983). From communicative competence to communicative language pedagogy. \textit{Language and Communication}, 1-47.

    \item Cassany, D. (2006). \textit{Tras las líneas: Sobre la lectura contemporánea}. Anagrama.

    \item Hattie, J., \& Timperley, H. (2007). The power of feedback. \textit{Review of Educational Research}, 77(1), 81-112.

    \item Hymes, D. (1972). On communicative competence. \textit{Sociolinguistics}, 269-293.

    \item Kapp, K. M. (2012). \textit{The Gamification of Learning and Instruction}. Pfeiffer.

    \item Marquès, P. (2010). El software educativo. \textit{DIM-UAB}.

    \item Ministerio de Educación y Formación Profesional. (2022). \textit{PISA 2022: Informe español}.

    \item Real Academia Española. (2010). \textit{Nueva gramática de la lengua española}. Espasa.

    \item Van Dijk, T. A. (1980). \textit{Texto y contexto: semántica y pragmática del discurso}. Cátedra.

    \item Vygotsky, L. S. (1978). \textit{Mind in society: The development of higher psychological processes}. Harvard University Press.
\end{enumerate}

\newpage
\section{Anexos}

\subsection{Anexo A: Modelo de pre-test / post-test (comprensión lectora)}

\textbf{Ejemplo de ítems (20 en total):}

\begin{enumerate}
    \item \textbf{Texto breve + preguntas de nivel literal, inferencial y crítico.}
    
    \item \textbf{Identifica el sujeto en la oración:}
    \begin{quote}
    "Ayer, mi hermano y yo fuimos al cine"
    \end{quote}
    \begin{enumerate}[label=\alph*)]
        \item Ayer
        \item Mi hermano y yo
        \item Al cine
        \item Fuimos
    \end{enumerate}
    
    \item \textbf{Señala la opción correcta (dequeísmo/queísmo):}
    \begin{enumerate}[label=\alph*)]
        \item Me dijo de que vendría
        \item Me dijo que vendría
        \item Me dijo en que vendría
    \end{enumerate}
\end{enumerate}

\subsection{Anexo B: Encuesta de percepción (Likert 1-5)}

\begin{center}
\begin{tabular}{|p{8cm}|c|c|c|c|c|}
\hline
\textbf{Afirmación} & \textbf{1} & \textbf{2} & \textbf{3} & \textbf{4} & \textbf{5} \\
\hline
Las mini lecciones HTML son fáciles de usar. & & & & & \\
\hline
La retroalimentación inmediata me ayudó a aprender. & & & & & \\
\hline
Prefiero practicar con ordenador que con papel. & & & & & \\
\hline
Los juegos gramaticales me motivaron. & & & & & \\
\hline
He mejorado mi escritura con estos recursos. & & & & & \\
\hline
Recomendaría estos recursos a otros compañeros. & & & & & \\
\hline
\end{tabular}
\end{center}

\subsection{Anexo C: Guía para grupo focal (45 min)}

\textbf{Preguntas guía:}

\begin{enumerate}
    \item ¿Qué fue lo que más te gustó de las lecciones interactivas? ¿Y lo que menos?
    \item ¿En qué aspecto de la lengua sientes que más has mejorado?
    \item ¿Cambiarías algo de los recursos? ¿El qué?
    \item ¿Preferirías seguir usando estas herramientas el próximo curso? ¿Por qué?
    \item ¿Cómo describirías la experiencia a un amigo que no la ha vivido?
\end{enumerate}

\subsection{Anexo D: Rúbrica de evaluación de texto argumentativo}

\begin{table}[H]
\centering
\caption{Rúbrica para evaluar textos producidos}
\label{tab:rubrica}
\begin{tabular}{|p{3cm}|p{2.5cm}|p{2.5cm}|p{2.5cm}|p{2.5cm}|}
\hline
\textbf{Criterio} & \textbf{Excelente (4)} & \textbf{Bueno (3)} & \textbf{Suficiente (2)} & \textbf{Insuficiente (1)} \\
\hline
Coherencia & Totalmente coherente & Mayormente coherente & Alguna incoherencia & Incoherente \\
\hline
Conectores & 5+ variados y precisos & 3-4 adecuados & 1-2 repetitivos & Ausentes \\
\hline
Ortografía & Sin errores & 1-2 errores & 3-4 errores & 5+ errores \\
\hline
Sintaxis & Oraciones complejas bien construidas & Algún error menor & Errores que afectan comprensión & Imposible de seguir \\
\hline
\end{tabular}
\end{table}

\subsection{Anexo E: Plan de lección tipo (50 minutos)}

\begin{table}[H]
\centering
\caption{Estructura de una sesión típica}
\label{tab:leccion}
\begin{tabular}{|p{2.5cm}|p{4cm}|p{5cm}|}
\hline
\textbf{Tiempo} & \textbf{Actividad} & \textbf{Descripción} \\
\hline
0-5' & Activación & Repaso rápido de la clase anterior \\
\hline
5-15' & Explicación & Introducción del nuevo concepto con mini HTML \\
\hline
15-40' & Práctica interactiva & Ejercicios con el programa o lección HTML \\
\hline
40-45' & Corrección & Resolución de dudas en grupo \\
\hline
45-50' & Cierre & Registro de progreso y avance \\
\hline
\end{tabular}
\end{table}

\end{document}
